\section{An abstract vanishing criterion for real-supported Euler data}
\label{sec:s5-crystal-criterion}

The positive-prime measure in the preceding argument does not need
an independent Tauberian growth estimate once the signed divisor is
known to be supported on the real height axis. The explicit-formula
axiom itself gives its polynomial growth. This yields the following
criterion for an actual $L$-datum; it is not restricted to the
$S_5$ comparison.

\paragraph{MA-FQ-006 (real integral divisor and eventual Euler positivity).}
Let $\mathcal F=(f,K,m)$ be an $L$-datum in the precise sense of
Booker, Definition~1.3. Suppose that
\[
 \begin{gathered}
 K(x)=0\quad(x>0),\qquad
 m(z)=0\quad(z\notin\mathbb R),\qquad
 m(t)\in\mathbb Z\quad(t\in\mathbb R),\\
 f(n)=0\quad\text{if }n\ge2\text{ is not a prime power},\\
 \text{there is a finite set of primes }\mathcal B
 \text{ such that }
 f(p^j)\in[0,\infty)
       \quad(p\notin\mathcal B,\ j\ge1).
 \end{gathered}
 \tag{MA-FQ.40}
\]
Then
\[
                  \mathcal F=(0,0,0).
 \tag{MA-FQ.41}
\]
In particular its associated $L$-function is identically one.
The hypothesis $K\equiv0$ is equality of the full kernel,
which is stronger than degree zero alone. No sign condition
on $m$, no assumption on its central value, and no separate
growth estimate for the good-prime measure are required.

\begin{proof}
\emph{The atomic divisor and the finite bad-prime correction.}
Use the ordinary Fourier convention (MA-FQ.1f):
\[
 \widehat u(x)=\int e^{itx}u(t)\,dt,\qquad
 \langle\widehat T,g\rangle=\langle T,\widehat g\rangle,
 \qquad
 u(t)=\frac1{2\pi}\int e^{-itx}\widehat u(x)\,dx.
 \tag{MA-FQ.42}
\]
By axiom (A3) there are constants $C,A$ with
\[
 \sum_{|t|\le T}|m(t)|\le C(1+T)^A.
 \tag{MA-FQ.43}
\]
Integrality makes the nonzero support locally finite in
the ordinary real topology. The signed atomic measure
$\mu=\sum_t m(t)\delta_t$ is tempered: summing (MA-FQ.43)
over dyadic annuli shows that its total variation is
integrable against $(1+|t|)^{-N}$ for $N>A$.

Define the positive locally finite measure
\[
 M_{\rm good}
   =\sum_{\substack{p\notin\mathcal B\\j\ge1}}
        f(p^j)\bigl(\delta_{j\log p}+\delta_{-j\log p}\bigr)
 \tag{MA-FQ.44}
\]
and the complex measure
\[
 B_0=2f(1)\delta_0+
    \sum_{\substack{p\in\mathcal B\\j\ge1}}
       \bigl(f(p^j)\delta_{j\log p}
                    +\overline{f(p^j)}\delta_{-j\log p}\bigr).
 \tag{MA-FQ.45}
\]
Local finiteness in (MA-FQ.44) follows because a bounded
interval contains logarithms of only finitely many
positive integers. Axiom (A1), with exponent two, gives
for every fixed prime $p$
\[
 \sum_{j\ge1}|f(p^j)|
       \le C_2\sum_{j\ge1}(j\log p)^{-2}<\infty.
 \tag{MA-FQ.46}
\]
Thus $B_0$ has finite total variation. This uses no sign
restriction at the bad primes. Also $f(1)\in\mathbb R$,
as required by (A1).

\emph{The explicit formula supplies the missing growth bound.}
For $g\in C_c^\infty$ with $g(-x)=\overline{g(x)}$,
its Fourier transform is real on $\mathbb R$.
Axiom (A4), with $K=0$, gives
\[
 \langle\mu,\widehat g\rangle
            =-2\operatorname{Re}
                    \sum_{n\ge1}f(n)g(\log n).
 \tag{MA-FQ.47}
\]
The sum includes $n=1$ under the minus sign. The
prime-power support and the signs at the good primes
give the identity
\[
                 \widehat{(-\mu)}=M_{\rm good}+B_0
                    \quad\text{on }C_c^\infty(\mathbb R).
 \tag{MA-FQ.48}
\]
This is a complex distribution identity, not merely an
identity on real tests. Indeed every $g$ decomposes as
$g=g_1+ig_2$, where
\[
 g_1(x)=\tfrac12\bigl(g(x)+\overline{g(-x)}\bigr),\qquad
 g_2(x)=\tfrac1{2i}\bigl(g(x)-\overline{g(-x)}\bigr).
 \tag{MA-FQ.49}
\]
Both have the requisite Hermitian symmetry, and both
sides of (MA-FQ.48) are complex linear in the test.

The left side of (MA-FQ.48) and $B_0$ are already
tempered. Hence the locally defined positive measure
$M_{\rm good}$ agrees on compact tests with the tempered
distribution $T=\widehat{(-\mu)}-B_0$.
Choose $\theta\in C_c^\infty(\mathbb R)$, nonnegative
and at least one on $[-1,1]$, and put
$\theta_R(x)=\theta(x/R)$ for $R\ge1$.
Continuity of $T$ on the Schwartz space bounds its
action by finitely many Schwartz seminorms. For some
integer $D$ and constant $C'$ this gives
\[
 \begin{aligned}
 M_{\rm good}([-R,R])
   &\le\langle M_{\rm good},\theta_R\rangle
      =\langle T,\theta_R\rangle\\
   &\le |\langle T,\theta_R\rangle|
      \le C'(1+R)^D.
 \end{aligned}
 \tag{MA-FQ.50}
\]
For clarity, every seminorm
$\sup_x|x^r\partial_x^\ell\theta_R(x)|$ is
$R^{r-\ell}\sup_y|y^r\theta^{(\ell)}(y)|$;
the finitely many such powers are bounded by one
fixed power of $1+R$. Positivity gives the first
inequality, and the middle pairing is real and
nonnegative. Dyadic summation now proves that
$M_{\rm good}$ is a tempered measure. Density of
$C_c^\infty$ in the Schwartz space extends
(MA-FQ.48) to an equality of tempered distributions.
No Tauberian theorem has entered this step.

\emph{Atomic positivity.}
Define the bounded continuous function and the tempered
distribution
\[
 q(t)=-\frac1{2\pi}\int e^{-itx}\,dB_0(x),
                 \qquad \nu=-\mu+q(t)\,dt.
 \tag{MA-FQ.51}
\]
Finite total variation of $B_0$ and Fourier inversion
give $\widehat{q(t)\,dt}=-B_0$, hence
$\widehat\nu=M_{\rm good}$.
For $\widetilde\psi(t)=\overline{\psi(-t)}$,
Fourier inversion and Fubini yield
\[
 \langle\nu,\psi*\widetilde\psi\rangle
      =\frac1{2\pi}\int
                    |\widehat\psi(-x)|^2\,dM_{\rm good}(x)
      \ge0\qquad(\psi\in C_c^\infty).
 \tag{MA-FQ.52}
\]
The integral is finite by the polynomial growth just proved.

We recall the precise atomic extraction so that it does not
depend on the $S_5$ specialization of MA-FQ-001.
Choose a real even $\phi\in C_c^\infty$ of $L^2$ norm one.
Let $\Psi=\phi*\widetilde\phi$, so $\Psi(0)=1$, and put
$\phi_\epsilon(t)=\epsilon^{-1/2}\phi(t/\epsilon)$.
Then $\phi_\epsilon*\widetilde\phi_\epsilon(t)
=\Psi(t/\epsilon)$.
For any finite $x_i\in\mathbb R$ and $c_i\in\mathbb C$,
substitute
$\psi_\epsilon=\sum_i c_i\phi_\epsilon(\,\cdot-x_i)$
in (MA-FQ.52).
Local finiteness of the atomic support isolates the atom
at each of the finitely many differences $x_i-x_j$
when $\epsilon$ is small. Each continuous-density term
has modulus at most
$\epsilon\|q\|_\infty\|\Psi\|_1$.
Letting $\epsilon$ tend to zero consequently gives
\[
 \sum_{i,j}c_i\overline{c_j}\,a(x_i-x_j)\ge0,
                     \qquad a(t)=-m(t).
 \tag{MA-FQ.53}
\]
Thus $a$ is positive definite on $\mathbb R_d$,
the additive real group with the discrete topology.
The two-by-two inequality gives
$|a(t)|\le a(0)=-m(0)$; this conclusion has not
assumed a sign or a prescribed value for $m(0)$.
If $a(0)=0$, the same inequality gives $a=0$.

\emph{Only finitely many primes can occur.}
The measure $\sum_t a(t)\delta_t=-\mu$ is tempered,
its coefficient function is integer valued and
positive definite on $\mathbb R_d$, and its
ordinary support is locally finite. Its ordinary
Fourier transform in (MA-FQ.48) is a locally finite
purely atomic measure. These are exactly the
hypotheses of MA-FQ-003. That theorem, proved from
Cohen's idempotent theorem, the real-line coset
classification and ordinary Poisson summation,
therefore gives a finite set of $h_\ell>0$ such that
\[
 \operatorname{supp}\widehat{(-\mu)}
                    \subseteq\bigcup_{\ell=1}^L h_\ell\mathbb Z.
 \tag{MA-FQ.54}
\]
The compact-dual representing measure on
$\widehat{\mathbb R_d}$ used in that proof is not
identified with this ordinary Fourier transform.
The empty union is permitted.

Call a good prime $p$ active when $f(p^j)>0$ for
at least one $j\ge1$. At $x=j\log p$ such a
coefficient is an atom of $\widehat{(-\mu)}$.
No other prime power has the same integer value
$p^j$, no bad-prime progression meets it, and the
conductor atom is at zero. Thus the atom cannot
cancel. Choose one such exponent $j_p$ for each
active prime. By (MA-FQ.54), $j_p\log p$ lies in
one of the finitely many frequency lattices.

Each one of those lattices can contain such
frequencies from at most one distinct prime,
even though the exponents $j_p$ may vary.
Indeed, if
$j_p\log p=nh$ and $j_q\log q=mh$ with
$m,n,j_p,j_q$ positive integers, then
$p^{mj_p}=q^{nj_q}$. Unique factorization forces
$p=q$. There are therefore only finitely many
active good primes.

\emph{Finite prime support forces the whole datum to vanish.}
Adjoin the finitely many bad primes to the active good
primes. The resulting finite set $\mathcal P$ contains
every prime whose powers support a nonzero coefficient
$f(n)$ for $n\ge2$. By (MA-FQ.46),
\[
                \sum_{n\ge2}|f(n)|
                =\sum_{p\in\mathcal P}\sum_{j\ge1}|f(p^j)|
                <\infty.
 \tag{MA-FQ.55}
\]
Consequently the complete Fourier measure
\[
 D=2f(1)\delta_0+
       \sum_{n\ge2}\bigl(f(n)\delta_{\log n}
                    +\overline{f(n)}\delta_{-\log n}\bigr)
          =\widehat{(-\mu)}
 \tag{MA-FQ.56}
\]
has finite total variation. Its inverse Fourier
transform is the bounded continuous density
\[
               r(t)=\frac1{2\pi}\int e^{-itx}\,dD(x).
 \tag{MA-FQ.57}
\]
Fourier inversion on Schwartz tests gives the
equality of tempered distributions $-\mu=r(t)\,dt$.
Both sides are locally finite measures.
The left side is purely atomic and the right
side absolutely continuous. Thus every atomic
mass is zero, so $\mu=0$ and $m=0$.
Equivalently uniqueness of the local Lebesgue
decomposition gives the same conclusion.
Taking Fourier transforms gives $D=0$.
The atom at zero is $2f(1)$, and the atom at
$\log n>0$ is exactly $f(n)$; these locations
are distinct. Hence $f(1)=0$ and $f(n)=0$
for every $n\ge2$. With $K=0$ this proves
(MA-FQ.41). The defining exponential series for
the associated $L$-function is then identically
one in $\operatorname{Re}s>1$.
\end{proof}

\paragraph{The finite nonintegral exceptions allowed by (A3).}
The same vanishing conclusion holds if the extra requirement
$m(t)\in\mathbb Z$ in (MA-FQ.40) is omitted and one retains
exactly Booker's axiom (A3), which allows only finitely many
points with nonintegral mass. Here is the complete modification.

First, the integral nonzero masses have absolute value at least
one, so (MA-FQ.43) bounds their number in every compact interval.
Adjoining the finitely many exceptional points preserves local
finiteness. The measure is tempered by the same dyadic estimate.
Every step through atomic positivity (MA-FQ.53) therefore applies
without change. Its two-by-two inequality bounds every value
of $a=-m$ by the finite real number $a(0)$. The integral values
in this bounded interval form a finite set, and the finitely
many exceptional points add only finitely many other values.
Thus $a$ has finite real range.

Bochner's theorem on $\mathbb R_d$, in the explicit
GNS--Riesz form proved in MA-FQ-001, gives
$a\in B(\mathbb R_d)$; that proof only uses positive
definiteness and the finite mass $a(0)$, not its integrality.
For each nonzero value $r$ in the finite range, choose a
polynomial taking value one at $r$ and zero at every other
range value. Applied to $a$, it gives
$1_{\{a=r\}}\in B(\mathbb R_d)$.
Cohen's idempotent theorem places this level set in the
coset ring. It is locally finite because the support of
$a$ is. The set assertion of MA-FQ-002, followed by
inclusion--exclusion as in that proof, expresses its
indicator as a finite integer linear combination of
ordinary lattice-coset and singleton indicators.
Multiplication by $r$ and summation over the finite
nonzero range expresses the atomic measure of $a$
as a finite \emph{real} linear combination of translated
lattice combs and singleton masses.

The ordinary Poisson formula (MA-FQ.24) applies to those
real coefficients just as it does to integer coefficients.
Its singleton terms give a continuous trigonometric-polynomial
density. Since the actual Fourier transform (MA-FQ.48) is
purely atomic, uniqueness of the local Lebesgue decomposition
eliminates that density. Its remaining support lies in
finitely many unshifted frequency lattices, exactly as in
(MA-FQ.54). The varying-prime-power argument, finite total
variation, inverse continuous density and atomwise vanishing
in the last two parts of the main proof now apply verbatim.
They give $\mathcal F=0$. This proves the extension using
precisely the finite-exception clause in (A3); the integral
case remains the stated main theorem.

\paragraph{Relation to the $S_5$ gate and to the scope tests.}
For the actual difference datum of MA-LD-001,
the gamma calculation in MA-FQ-001 gives $K_R=0$,
the signed orders are integral, and its good
prime-power coefficients are nonnegative.
The required real support has a shorter direct
proof from MT-S508, without MA-S5-004.
In its paragraph \emph{The positive Euler quotient
controls a whole half-plane}, the proof following
(S508.13) uses (S508.9), the positive Dirichlet
series and Landau's theorem.  Only the sufficient
conclusion $\gamma\le1/2$ is needed here.
Indeed the source first gives $\gamma<1$.
If $1/2<\gamma<1$, Landau's singularity at the real
abscissa would be a pole of the meromorphic quotient
at a point where $(\mathcal R_B)$ makes its
denominator nonzero, a contradiction.
Any case $\gamma<1/2$, including $\gamma\le0$,
already gives the required half-plane convergence;
no equality assertion about that abscissa is needed.
A finite $\Omega$ also causes no convergence problem.
For every $\sigma>1/2$, the convergent positive
Dirichlet series has prime coefficients equal to
four, so $\sum_{p\in\Omega}p^{-\sigma}<\infty$.
Its Euler logarithm therefore converges absolutely
and locally uniformly in $\operatorname{Re}s>1/2$.
The finitely many bad-prime logarithms converge
there by (A1), as in MA-LD-001. Thus the actual
finite quotient $A/L_W$ is holomorphic and
nonvanishing throughout that half-plane.

The exactly equal gamma factors cancel, and
the conductor exponential is nowhere zero.
Consequently the completed quotient also has
no divisor on $\operatorname{Re}s>1/2$,
including at $s=1$. Its completed functional
equation reflects the full divisor across
$\operatorname{Re}s=1/2$, so there is no divisor
on the open left half-plane either, including
at $s=0$. This verifies real height support
with no omitted endpoint poles. MA-FQ-006 now
applies directly. Its vanishing includes every
bad-prime logarithmic coefficient and the
conductor atom; equality of the local logarithmic
series implies equality of the normalized
standard Euler factors at every finite prime.
The archimedean factors were already equal.
It also gives $m_R(0)=-k=0$.

Thus MA-S5-004 and the original MA-LD-001
proof remain valid historical deductions, but
their separate Tauberian exceptional-prime
estimate is not needed for this final closure.
Temperedness of the positive prime measure
comes directly from (MA-FQ.50).
The hypothesis $(\mathcal R_B)$ remains an
explicit premise and is not proved here.

The exact hypotheses explain the comparisons
in MA-FQ-005. The ratio $\zeta/L(\,\cdot,\chi_5)$
has completed poles off the critical line;
the equal-parity Dirichlet ratio there has
infinitely many negative good-prime coefficients;
the single-prime central Euler pole violates
(A1); and the stated nonperiodic positive
Fourier crystal has mixed logarithmic
frequencies outside the prime-power support.
Those checked examples therefore do not meet
(MA-FQ.40).

\paragraph{Primary inputs and attribution.}
The exact axioms used are Booker (2015),
\emph{L-functions as distributions},
Definition~1.3(A1), (A3), (A4), author p.~3
of \path{arXiv:1308.3067v2}.
The finite-prime-support conclusion also
satisfies the hypothesis of Booker's
Theorem~1.6(ii), author p.~5, because
(MA-FQ.55) implies
$\sum_{n\ge2}|f(n)|/\log n<\infty$.
The last part of the proof above supplies
the needed vanishing directly in the
real-supported, zero-kernel case.
The deep harmonic-analysis input remains
Cohen's idempotent theorem in the precise
form checked in MA-FQ-002--003:
Sanders (2020), \S12, Theorems~12.1 and~12.3,
and the real-line locally finite coset
argument supplied in MA-FQ-002, with
Schreiber (1970), Lemma~1.4 and
Corollary~1.5, p.~807, Example~1.6,
p.~808, as its exact primary locator.
The criterion is a consequence of those
inputs and the explicit-formula axioms;
no claim of priority over known rigidity
results is made.

Booker, A. R. (2015). $L$-functions as distributions.
\emph{Mathematische Annalen, 363}, 423--454.
\url{https://doi.org/10.1007/s00208-015-1178-z}.

Sanders, T. (2020). Bounds in Cohen's idempotent theorem.
\emph{Journal of Fourier Analysis and Applications, 26},
Article~25.
\url{https://doi.org/10.1007/s00041-020-09732-y}.

Schreiber, B. M. (1970). On the coset ring and strong Ditkin sets.
\emph{Pacific Journal of Mathematics, 32}(3), 805--812.
\url{https://doi.org/10.2140/pjm.1970.32.805}.
